\section{Introdução}
\label{sec:introducao}

Impulsionada pela Mineração de Dados Educacionais (MDE), a tomada de decisão orientada a dados vem sendo adotada na gestão educacional para alocar recursos, prever a evasão e estimar o desempenho acadêmico de estudantes \cite{romero2010educational}. No Brasil, os microdados do Exame Nacional do Ensino Médio (ENEM), publicados pelo Instituto Nacional de Estudos e Pesquisas Educacionais Anísio Teixeira (INEP) \cite{inep2022microdados}, ocupam posição central nessas iniciativas \cite{luz2025uso, souza2025desigualdades}: por ser o principal mecanismo de acesso ao ensino superior no país, o exame alimenta modelos que buscam estimar o sucesso acadêmico ou direcionar políticas de assistência.

Modelos de alta capacidade preditiva, contudo, operam frequentemente como ``caixas-pretas'', cujos caminhos de decisão não são inteligíveis para seres humanos, o que dificulta a responsabilização e a contestação de decisões \cite{adadi2018peeking, rudin2019stop}. Além disso, esses modelos herdam e podem amplificar os vieses históricos dos dados de treinamento \cite{barocas2023fairness}. Aplicados a bases que registram desigualdades estruturais, como diferenças de renda familiar, de acesso a tecnologia e de raça/cor, tendem a convertê-las em regras automatizadas que desfavorecem grupos historicamente vulneráveis, sob a aparência de objetividade técnica \cite{baker2022algorithmic}.

Nesse cenário, a Inteligência Artificial Explicável (XAI, do inglês \textit{Explainable Artificial Intelligence}) e a Justiça Algorítmica (\textit{Algorithmic Fairness}) assumem relevância ética e regulatória \cite{arrieta2020explainable, barocas2023fairness}. No Brasil, a Lei Geral de Proteção de Dados Pessoais (LGPD) assegura ao titular a revisão de decisões tomadas unicamente com base em tratamento automatizado \cite{brasil2018lgpd}, e o Projeto de Lei nº 2.338/2023 propõe avaliações de impacto algorítmico para sistemas de alto risco \cite{brasil2023pl2338}. As bibliotecas existentes, contudo, exigem integração programática especializada \cite{bird2020fairlearn, bellamy2018aif360}, e faltam ferramentas de auditoria sistemática e reprodutível acessíveis a gestores e pesquisadores sem formação em programação.

Este artigo busca responder à seguinte questão de pesquisa: \textit{é possível construir uma ferramenta de auditoria, executável em infraestrutura gratuita, que detecte, quantifique e explique vieses de raça/cor e renda em modelos preditivos do ENEM, e que seja sensível a vieses aos quais as métricas convencionais de desempenho pouco respondem?} O objetivo é projetar, desenvolver e validar essa ferramenta, denominada AuditEdu-XAI, que integra o \textit{framework} SHAP \cite{lundberg2017unified} a métricas de Impacto Disparate\footnote{Adota-se ``Impacto Disparate'' como decalque do termo \textit{Disparate Impact}, consagrado na literatura de \textit{fairness} \cite{feldman2015certifying}; a expressão designa a razão entre taxas de seleção, e não o sentido corrente da palavra ``disparate'' em português.} \cite{feldman2015certifying} e de Igualdade de Oportunidades \cite{hardt2016equality}, segundo a metodologia de \textit{Design Science Research} (DSR) \cite{peffers2007design}.

As contribuições do trabalho são:
\begin{itemize}
    \item um artefato de software de arquitetura desacoplada (pipeline em lote e interface que apenas lê resultados pré-computados), com código publicado em repositório público, \url{https://github.com/rbrtotavio/audit_edu_xai}, implantado gratuitamente em \url{https://auditeduxai.streamlit.app/} e empacotado em contêiner Docker validado por integração contínua;
    \item a auditoria de classificadores XGBoost treinados com as edições de 2019, 2020 e 2021 do ENEM, restrita ao conjunto de teste, com intervalos de confiança por \textit{bootstrap} e a decomposição do viés entre dados e modelo;
    \item um protocolo de validação por injeção sintética de viés com controle negativo pareado, que separa o efeito do viés direcionado do efeito genérico de alterar a prevalência de rótulos.
\end{itemize}

O restante do artigo está organizado da seguinte forma. A Seção~\ref{sec:fundamentacao} apresenta a fundamentação teórica e os trabalhos relacionados. A Seção~\ref{sec:metodologia} descreve a metodologia, e a Seção~\ref{sec:artefato}, o artefato. A Seção~\ref{sec:resultados} discute os resultados, e a Seção~\ref{sec:conclusao} conclui.
